% SGA 3, 第 VI_B 篇, 第 1 轮英文重构.
% 范围: 命题 4.2, 含底本编者注 28--30.
% 底本: Polo--Gille Exp6B-13oct24.pdf, 分篇第 22--24 页
% (合订本第 363--365 页; 原版页边印刷页码 348--350).
% OCR 与 Reinhold 英文 Markdown 仅为比对材料.

\medskip
\phantomsection\label{prop:sga3-VIB-4.2}%
\noindent\textbf{命题 4.2.---}
\emph{设 \(\pi:G\to S\) 是局部有限呈示的 \(S\)-概形,
带有一个 \(S\)-截面 \(\epsilon\), 且满足以下两个条件
\textup{(当 \(G\) 是 \(S\)-群时, 由 \hyperref[lem:sga3-VIB-1.5]{1.5}
及 VI\(_A\), \SGARefLink{cor:sga3-VIA-2.4.1}{2.4.1}, 这两个条件成立)}:}
\begin{enumerate}
\item[\textup{\SGAInlineRefTarget{prop:sga3-VIB-4.2:item:a}(a)}]
  \emph{对每个 \(s\in S\) 及每个 \(x\in G_s\), 有}
  \[
    \dim G_s=\dim_x G_s
  \]
  \emph{\textup{(等价地 \textup{(\hyperref[lem:sga3-VIB-1.5]{1.5})},
  对每个 \(s\in S\), \(G_s\) 的所有不可约分支维数相同)}.}
\item[\textup{(b)}] \emph{对每个 \(s\in S\), 若用
  \(\underline{G_s^0}\) 表示 \(G_s\) 中包含 \(\epsilon(s)\)
  的连通分支, 则 \(\underline{G_s^0}\) 几何不可约.}
\end{enumerate}
\emph{设 \(s\in S\). 下列条件等价:}
\begin{enumerate}
\item[\textup{\SGAInlineRefTarget{prop:sga3-VIB-4.2:item:i}(i)}]
  \emph{\(G\) 在 \(\underline{G_s^0}\) 的各点处对 \(S\) 普遍开.}
\item[\textup{(i bis)}] \emph{\(G\) 在 \(\epsilon(s)\)
  在 \(\underline{G_s^0}\) 中的某邻域的每个点处对 \(S\) 普遍开.}%
  \footnote{\label{note:sga3-VIB-28}编者注: 原文作``在 \(\epsilon(s)\) 的各点处''.
  仅假设 \(G\) 在 \(\epsilon(s)\) 处对 \(S\) 普遍开是不够的,
  如编者\hyperref[note:sga3-VIB-19]{注 19} 中的例子所示.
  我们已相应地改正证明.}
\item[\textup{\SGAInlineRefTarget{prop:sga3-VIB-4.2:item:ii}(ii)}] \emph{函数}
  \[
    t\longmapsto\dim G_t
  \]
  \emph{在 \(s\) 在 \(S\) 中的某邻域上为常数.}
\item[\textup{(iii)}]\phantomsection\label{concl:sga3-VIB-4.2-iii}
  \emph{\(\underline{G^0}\) ``在 \(\underline{G_s^0}\) 的各点处对 \(S\) 普遍开'',
  即给定 \(S'\to S\), \(s'\in {S'}_s\),
  \(g\in\underline{G_{s'}^0}\), 以及 \(g\) 在 \(G'=G_{S'}\)
  中的开邻域 \(V\), 则}
  \[
    \pi\bigl(V\cap\underline{{G'}^0}\bigr)
  \]
  \emph{是 \(s'\) 在 \(S'\) 中的开邻域.}%
  \footnote{\label{note:sga3-VIB-29}编者注: 新增了由 O.\ Gabber 指出的条件
  \textup{\SGARefLink{concl:sga3-VIB-4.2-iii}{(iii)}},
  后文将用到它 \textup{(\SGARefLink{thm:sga3-VIB-5.7}{5.7})}.}
\end{enumerate}

显然 (i) 蕴含 (i bis). 由 EGA~IV\(_3\), 14.3.3.1(ii),
\(\underline{G_s^0}\) 中使 \(\pi_G\) 普遍开的点所成的集合
在 \(\underline{G_s^0}\) 中闭. 由于 \(\underline{G_s^0}\) 不可约,
可知 (i bis) 蕴含 \SGARefLink{prop:sga3-VIB-4.2:item:i}{(i)}.

证明 \SGARefLink{prop:sga3-VIB-4.2:item:i}{(i)} 蕴含
\SGARefLink{prop:sga3-VIB-4.2:item:ii}{(ii)}.
设 \(T\) 是满足
\[
  \dim G_t=\dim G_s.
\]
的 \(t\in S\) 所成的集合. 由 4.1, \(T\) 局部可构造;
因此由 EGA~IV\(_1\), 1.10.1, 为证明 \(T\) 是 \(s\) 的邻域,
只须证明 \(s\) 的每个泛化 \(s'\) 都属于 \(T\).

设 \(x\) 是 \(\underline{G_s^0}\) 的泛点,
\(U\) 是 \(G\) 的包含 \(x\) 的开集.
由于 \(\pi_G\) 在 \(x\) 处普遍开,
EGA~IV\(_3\), 14.3.13 对 \(s\) 的每个泛化 \(s'\) 给出
\[
  \dim\bigl(U\cap G_{s'}\bigr)
    \geq \dim_x\bigl(U\cap G_s\bigr).
\]
由假设 \SGARefLink{prop:sga3-VIB-4.2:item:a}{(a)}, 这蕴含
\[
  \dim G_{s'}\geq\dim G_s.
\]
另一方面, 由 \hyperref[prop:sga3-VIB-4.1]{4.1},
函数 \(s\mapsto\dim G_s\) 上半连续, 因此
\[
  \dim G_{s'}\leq\dim G_s.
\]
于是 \(s'\in T\), 证明了 \SGARefLink{prop:sga3-VIB-4.1:item:i}{(i)}
蕴含 \SGARefLink{prop:sga3-VIB-4.1:item:ii}{(ii)}.

显然 \SGARefLink{concl:sga3-VIB-4.2-iii}{(iii)} 蕴含
\SGARefLink{prop:sga3-VIB-4.1:item:i}{(i)};
证明 \SGARefLink{prop:sga3-VIB-4.2:item:ii}{(ii)} 蕴含
\SGARefLink{concl:sga3-VIB-4.2-iii}{(iii)}.
纤维维数在底域扩张下不变, 而 \(\underline{G^0}\) 的构造
与基变换可交换 (见 \hyperref[prop:sga3-VIB-3.3]{3.3} 的证明).
因此可假设 \(S'=S\) 且 \(s'=s\).
此外, \(G\) 中每个与 \(\underline{G_s^0}\) 相交的开集 \(V\)
都包含 \(\underline{G_s^0}\) 的泛点 \(\eta\), 故可假设 \(g=\eta\).

还可假设 \(S\) 仿射. 设 \(U\) 是 \(\epsilon(s)\) 的仿射开邻域;
则它在 \(S\) 上有限呈示. 将 \(S\) 改为 \(\epsilon^{-1}(U)\),
再将 \(U\) 改为 \(U\cap\pi^{-1}(S)\), 即归结为
\[
  \pi:U\longrightarrow S
\]
有限呈示且容许 \(\epsilon:S\to U\) 为截面的情形.
由 EGA~IV\(_3\), 9.7.12, \(\underline{G^0}\cap U\)
在 \(U\) 中可构造. 将 \(V\) 改为包含在 \(V\cap U\) 中的仿射开集,
可知 \(\underline{G^0}\cap V\) 在 \(V\) 中可构造.
由于 \(\pi:V\to S\) 有限呈示,
Chevalley 可构造性定理 (EGA~IV\(_1\), 1.8.4) 表明
\[
  \pi\bigl(\underline{G^0}\cap V\bigr)
\]
在 \(S\) 中局部可构造.

因此, 由\emph{同上}, 1.10.1, 为证明
\(\pi(\underline{G^0}\cap V)\) 是 \(s\) 的开邻域,
只须证明, 对 \(s\) 的每个泛化 \(t\), 存在 \(\eta\)
的一个属于 \(\underline{G^0}\) 的泛化 \(\xi\),
从而它属于 \(\underline{G^0}\cap V\).
\(\underline{G_t^0}\) 的泛点 \(\xi\) 就是 \(\eta\) 的这样一个泛化.
事实上, \(\epsilon(s)\) 属于 \(\{\xi\}\) 在 \(G\) 中的闭包 \(X\);
因而由 Chevalley 半连续性定理 (EGA~IV\(_3\), 13.1.3),
\[
  \dim_{\epsilon(s)}X_s\geq\dim_\xi X_t.
\]
另一方面, 假设 \SGARefLink{prop:sga3-VIB-4.2:item:ii}{(ii)} 蕴含
\[
  \dim_\xi\underline{G_t^0}=\dim_{\epsilon(s)}G_s.
\]
因此 \(X_s\) 中包含 \(\epsilon(s)\) 的某个不可约分支等于
\(\underline{G_s^0}\), 从而
\[
  \eta\in\overline{\{\xi\}}.
\]
这证明了 \SGARefLink{prop:sga3-VIB-4.2:item:ii}{(ii)} 蕴含
\SGARefLink{concl:sga3-VIB-4.2-iii}{(iii)}, 因而证明了命题.

\footnote{\label{note:sga3-VIB-30}编者注: 前面的论证由 O.\ Gabber 告知我们.
我们也保留了原文中给出的 \textup{\SGARefLink{prop:sga3-VIB-4.2:item:ii}{(ii)}}
\(\Rightarrow\)\textup{\SGARefLink{prop:sga3-VIB-4.2:item:i}{(i)}} 的证明.}%
还可如下证明 \SGARefLink{prop:sga3-VIB-4.2:item:ii}{(ii)} 蕴含
\SGARefLink{prop:sga3-VIB-4.2:item:i}{(i)}.
由于 \(\pi\) 局部有限呈示, 存在 \(G\) 的包含 \(\epsilon(s)\) 的开集 \(U\),
以及 \(S\) 的包含 \(s\) 的开集 \(V\),
使 \(\pi(U)\subset V\), 且由 \(\pi\) 诱导的态射
\[
  \pi':U\longrightarrow V
\]
有限呈示. 置
\[
  T=\epsilon^{-1}(U),
  \qquad
  W={\pi'}^{-1}(T)=U\cap\pi^{-1}(T).
\]
由 \(\pi'\) 诱导的态射
\[
  \pi'':W\longrightarrow T
\]
有限呈示, 并容许由 \(\epsilon\) 诱导的态射
\[
  \epsilon'':T\longrightarrow W
\]
为截面. 此外, 对每个 \(t\in T\), 由于 \(\underline{G_t^0}\) 不可约,
\(W\cap\underline{G_t^0}\) 在 \(\underline{G_t^0}\) 中稠密,
因而不可约, 从而连通. 所以它是 \({\pi''}^{-1}(t)\)
中包含 \(\epsilon''(t)\) 的连通分支.

由于 \(W\cap\underline{G_t^0}\) 是 \(\underline{G_t^0}\) 的稠密开集,
\hyperref[lem:sga3-VIB-1.5]{1.5} 及 EGA~IV\(_2\), 5.2.1 给出
\[
  \dim(W\cap\underline{G_t^0})=\dim\underline{G_t^0}=\dim G_t.
\]
于是函数
\[
  t\longmapsto\dim(W\cap\underline{G_t^0})
\]
在 \(s\) 在 \(T\) 中的某邻域上为常数.
最后证明对每个 \(t\in T\), \(W\cap\underline{G_t^0}\) 几何不可约.
设 \(K\) 是 \(\kappa(t)\) 的扩张. 则
\[
  (W\cap\underline{G_t^0})\otimes_{\kappa(t)}K
    =
  \bigl(W\otimes_{\kappa(t)}K\bigr)
    \cap
  \bigl(\underline{G_t^0}\otimes_{\kappa(t)}K\bigr)
\]
是 \(\underline{G_t^0}\otimes_{\kappa(t)}K\) 的非空开集,
因而不可约, 因为 \(\underline{G_t^0}\otimes_{\kappa(t)}K\) 不可约.

现在处于 EGA~IV\(_3\), 15.6.6(ii) 的情形,
该结果断言 \(\pi''\), 从而 \(\pi\),
在 \(W\cap\underline{G_s^0}\) 的各点处普遍开.
但由 EGA~IV\(_3\), 14.3.3.1(ii),
\(G_s\) 中使 \(\pi\) 普遍开的点所成的子集在 \(G_s\) 中闭.
由于它包含 \(W\cap\underline{G_s^0}\), 故也包含其闭包 \(\underline{G_s^0}\).
